\documentclass[11pt,a4paper]{article}
\usepackage[margin=2.5cm]{geometry}
\usepackage[T1]{fontenc}
\usepackage{amsmath,amssymb,amsthm,mathtools}
\usepackage{enumitem}
\usepackage{microtype}
\usepackage{xcolor}
\usepackage{hyperref}
\hypersetup{colorlinks=true,linkcolor=blue!50!black,citecolor=blue!50!black}

\newtheorem{theorem}{Theorem}
\newtheorem{lemma}[theorem]{Lemma}
\newtheorem{corollary}[theorem]{Corollary}
\theoremstyle{definition}
\newtheorem{definition}[theorem]{Definition}
\theoremstyle{remark}
\newtheorem{remark}[theorem]{Remark}

\title{Edge-disjoint barbells from edge-disjoint cycles}
\date{}

\begin{document}
\maketitle

\begin{abstract}
Let \(G\) be a connected multigraph that contains \(k\) edge-disjoint cycles.
We show that \(G\) contains at least \(\lceil (k-1)/3\rceil\) edge-disjoint
barbells, where a barbell consists of two edge-disjoint cycles joined by a
path. This bound needs no assumption on the lengths of the cycles, and it is
exact for infinitely many \(k\). As a corollary, if the cycles have length at
most \(\alpha\), with \(\alpha\ge2\) and \(k\ge2\), then there are at least
\(k/(2\alpha)\) edge-disjoint barbells.
\end{abstract}

\section{Definitions and statement}

Graphs are finite multigraphs; parallel edges and self-loops are allowed. A
\emph{cycle} of length \(L\ge1\) is a self-loop (\(L=1\)), a pair of parallel
edges (\(L=2\)), or, for \(L\ge3\), a closed path through \(L\) distinct
vertices. A \emph{path} may consist of a single vertex and no edges. For a
subgraph \(H\) we write \(V(H)\) and \(E(H)\) for its vertex and edge sets.

\begin{definition}[Barbell]\label{def:barbell}
A \emph{barbell} in \(G\) is a triple \(B=(D_1,D_2,P)\), where \(D_1,D_2\) are
edge-disjoint cycles of \(G\) and \(P\) is a path in \(G\) from a vertex of
\(D_1\) to a vertex of \(D_2\) such that
\begin{enumerate}[label=(\alph*)]
\item \(E(P)\cap\bigl(E(D_1)\cup E(D_2)\bigr)=\emptyset\), and
\item no interior vertex of \(P\) lies in \(V(D_1)\cup V(D_2)\).
\end{enumerate}
The cycles \(D_1,D_2\) are the \emph{ends} of \(B\), and
\(E(B):=E(D_1)\cup E(D_2)\cup E(P)\). Barbells \(B_1,\dots,B_b\) are
\emph{edge-disjoint} if the sets \(E(B_i)\) are pairwise disjoint.
\end{definition}

If \(D_1\) and \(D_2\) share a vertex, \(P\) may be that single vertex
(a ``figure eight''). Condition~(b) is not needed for anything below; we
include it only because it makes the conclusion stronger.

\begin{theorem}\label{thm:main}
Let \(G\) be a connected multigraph that contains \(k\) pairwise edge-disjoint
cycles \(C_1,\dots,C_k\). Then \(G\) contains at least
\[
  \left\lceil\frac{k-1}{3}\right\rceil
\]
edge-disjoint barbells, and every end of each of them is one of the cycles
\(C_i\). If every \(C_i\) has length at most \(2\), then \(G\) contains at least
\(\lfloor k/2\rfloor\) edge-disjoint barbells.
\end{theorem}

\begin{corollary}\label{cor:alpha}
Let \(G\) be a connected multigraph that contains \(k\ge2\) edge-disjoint
cycles, each of length at most \(\alpha\), where \(\alpha\ge2\). Then \(G\)
contains at least \(k/(2\alpha)\) edge-disjoint barbells. Without the
assumptions \(k\ge2\) and \(\alpha\ge2\), it still contains at least
\(\lfloor k/(2\alpha)\rfloor\) of them.
\end{corollary}

\begin{proof}
Let \(b\) be the number of barbells given by Theorem~\ref{thm:main}. If
\(k\in\{2,3\}\), then \(b\ge\lceil(k-1)/3\rceil=1\ge3/4\ge k/(2\alpha)\). If
\(k\ge4\), then \(b\ge(k-1)/3\ge k/4\ge k/(2\alpha)\). For the second claim,
if \(\alpha\ge2\) the bound just proved applies whenever \(k\ge2\), and for
\(k\le1\) we have \(\lfloor k/(2\alpha)\rfloor=0\). If \(\alpha=1\), every
\(C_i\) is a self-loop, and Theorem~\ref{thm:main} gives
\(b\ge\lfloor k/2\rfloor=\lfloor k/(2\alpha)\rfloor\).
\end{proof}

\begin{remark}[The assumptions in Corollary~\ref{cor:alpha} are needed]
\label{rem:assumptions}
The real-valued bound \(k/(2\alpha)\) can fail if \(k=1\) or \(\alpha=1\). A
single triangle has \(k=1\) but no barbell, although \(1/(2\alpha)>0\). A path
on three vertices with a self-loop at each vertex has \(k=3\) and
\(\alpha=1\), and it has only three cycles in total, so it contains at most one
edge-disjoint barbell, which is less than \(k/(2\alpha)=3/2\).
\end{remark}

\section{Two lemmas}

\begin{lemma}[From walks to barbells]\label{lem:walk}
Let \(D_1,D_2\) be edge-disjoint cycles of \(G\), and let \(W\) be a walk from a
vertex of \(D_1\) to a vertex of \(D_2\) with
\(E(W)\cap(E(D_1)\cup E(D_2))=\emptyset\). Then there is a path \(P\) with
\(E(P)\subseteq E(W)\) such that \((D_1,D_2,P)\) is a barbell.
\end{lemma}

\begin{proof}
If \(D_1\) and \(D_2\) share a vertex \(x\), take \(P\) to be the path
consisting of \(x\) alone. Otherwise, every walk contains a path between its
endpoints whose edges are edges of the walk, so let \(Q=q_0q_1\cdots q_s\) be
such a path inside \(W\), with \(q_0\in V(D_1)\) and \(q_s\in V(D_2)\). Let
\(j\) be the smallest index with \(q_j\in V(D_2)\), and let \(i\) be the
largest index \(i\le j\) with \(q_i\in V(D_1)\); it exists because
\(q_0\in V(D_1)\), and \(i<j\) because no vertex lies in both cycles. The
subpath \(P=q_i\cdots q_j\) satisfies condition~(a) because
\(E(P)\subseteq E(W)\), and condition~(b) by the choice of \(i\) and \(j\).
\end{proof}

\begin{lemma}[Disjoint arcs on a cycle]\label{lem:arcs}
Let \(C\) be a cycle with distinct vertices \(v_0,v_1,\dots,v_{L-1}\) and edges
\(v_tv_{t+1}\) (indices modulo \(L\)), and let
\(0<i_1<i_2<\cdots<i_r<L\). For \(1\le j<r\) let \(A_j\) be the path
\(v_{i_j}v_{i_j+1}\cdots v_{i_{j+1}}\), and let \(A_r\) be the path
\(v_{i_r}v_{i_r+1}\cdots v_{L-1}v_0\). Then the arcs \(A_1,A_3,A_5,\dots\) and,
if \(r\) is odd, \(A_r\), are pairwise edge-disjoint paths in \(C\).
\end{lemma}

\begin{proof}
The arc \(A_j\) uses exactly the edges \(v_tv_{t+1}\) with
\(t\in[i_j,i_{j+1})\), where \(i_{r+1}:=L\). For the listed indices \(j\)
these intervals are pairwise disjoint.
\end{proof}

\section{Proof of Theorem~\ref{thm:main}}

\paragraph{An auxiliary tree.}
Call an edge of \(G\) a \emph{cycle edge} if it belongs to some \(C_i\), and a
\emph{free edge} otherwise. Let \(A\) be the graph with vertex set
\(V(G)\cup\{c_1,\dots,c_k\}\), where each \(c_i\) is a new \emph{cycle node},
and with the following edges:
\begin{itemize}
\item every free edge of \(G\), with the same endpoints;
\item for every \(i\) and every vertex \(v\in V(C_i)\), one \emph{star edge}
  \(c_iv\).
\end{itemize}
The graph \(A\) is connected. Indeed, every edge \(xy\) of \(G\) is either a
free edge of \(A\), or an edge of some \(C_i\), and then \(x,y\in V(C_i)\) are
joined by the path \(x\,c_i\,y\) in \(A\). Since \(G\) is connected, all of
\(V(G)\) lies in one component of \(A\), and every \(c_i\) is adjacent to a
vertex of \(G\).

Fix a spanning tree \(T\) of \(A\) and root it at a vertex \(\rho\in V(G)\).
For a node \(x\) of \(T\), let \(T_x\) be the subtree of its descendants,
including \(x\). The neighbors of a cycle node \(c_i\) in \(A\) are vertices of
\(C_i\), and there is only one star edge for each of them. Hence \(c_i\) has a
parent \(p_i\in V(C_i)\), and its children are distinct vertices of
\(C_i\) different from \(p_i\). The tree edges between two vertices of \(G\) are
free edges, and a self-loop is never a tree edge.

For a node \(x\), define its \emph{resources} \(R_x\subseteq E(G)\) as the set
of free edges that are tree edges of \(T_x\), together with all edges of the
cycles \(C_i\) with \(c_i\in T_x\). If \(x\) and \(y\) are distinct children of
the same node, then \(T_x\) and \(T_y\) are disjoint, and since the cycles
\(C_i\) are edge-disjoint and free edges lie in no cycle,
\(R_x\cap R_y=\emptyset\).

\paragraph{Packages.}
A \emph{package at a vertex \(v\in V(G)\)} is a pair \((C,W)\), where \(C\) is
one of the cycles \(C_i\) and \(W\) is a walk in \(G\) from a vertex of \(C\) to
\(v\) with \(E(W)\cap E(C)=\emptyset\). We write \(E(C,W):=E(C)\cup E(W)\).

\paragraph{The inductive claim.}
For a node \(x\) of \(T\), let \(\hat x:=x\) if \(x\in V(G)\), and
\(\hat x:=p_i\) if \(x=c_i\). We prove by induction from the leaves of \(T\) to
the root that every node \(x\) admits
\begin{itemize}
\item a family \(\mathcal B_x\) of barbells,
\item a set \(\mathcal L_x\) of cycles (the \emph{lost} cycles), and
\item either nothing or one package \(\pi_x=(C^{x},W^{x})\) at \(\hat x\) (the
  \emph{exported} package),
\end{itemize}
with the following properties.
\begin{enumerate}[label=(P\arabic*)]
\item\label{P:partition} Every cycle \(C_i\) with \(c_i\in T_x\) is an end of
  exactly one barbell of \(\mathcal B_x\), or lies in \(\mathcal L_x\), or is
  \(C^{x}\); these three cases are mutually exclusive. All ends of barbells in
  \(\mathcal B_x\) are cycles \(C_i\) with \(c_i\in T_x\).
\item\label{P:disjoint} The sets \(E(B)\) for \(B\in\mathcal B_x\), together
  with \(E(\pi_x)\) if there is an exported package, are pairwise disjoint
  subsets of \(R_x\).
\item\label{P:count} \(|\mathcal L_x|\le|\mathcal B_x|\).
\end{enumerate}

\paragraph{Collecting the packages of the children.}
Let \(x\) be a node with children \(y_1,\dots,y_s\) (possibly \(s=0\)), and
assume that the claim holds for every child. We turn the exported packages of
the children into packages at a vertex of \(G\):
\begin{itemize}
\item If \(x\in V(G)\) and \(y_j\in V(G)\), the tree edge \(e_j=y_jx\) is a
  free edge. We replace \(\pi_{y_j}=(C,W)\), a package at \(y_j\), by
  \((C,We_j)\), a package at \(x\). Since \(e_j\) is free, it is not in
  \(E(C)\).
\item If \(x\in V(G)\) and \(y_j=c_i\) is a cycle node, then \(p_i=x\), so
  \(\pi_{y_j}\) is already a package at \(x\).
\item If \(x=c_i\) is a cycle node, its child \(y_j\) is a vertex of \(C_i\),
  and \(\pi_{y_j}\) stays a package at \(y_j\).
\end{itemize}
Call the resulting packages the \emph{arrivals} at \(x\). In each case the edge
set of an arrival from \(y_j\) is contained in \(R_{y_j}\cup\{e_j\}\), and
\(e_j\) belongs to \(R_x\) but to no \(R_{y_{j'}}\). Hence, by~\ref{P:disjoint}
for the children, all barbells in \(\bigcup_j\mathcal B_{y_j}\) and all arrivals
have pairwise disjoint edge sets, and these are contained in \(R_x\). Moreover,
for an arrival \((C,W)\) from \(y_j\) and a cycle \(C'\) belonging to an arrival
from \(y_{j'}\) with \(j'\ne j\), we have \(E(W)\cap E(C')=\emptyset\).

We start from
\(\mathcal B_x:=\bigcup_j\mathcal B_{y_j}\) and
\(\mathcal L_x:=\bigcup_j\mathcal L_{y_j}\), and add to them as follows.

\paragraph{Case 1: \(x=v\) is a vertex of \(G\).}
All arrivals \((C^{(1)},W^{(1)}),\dots,(C^{(r)},W^{(r)})\) are packages at
\(v\). For \(t=1,2,\dots,\lfloor r/2\rfloor\), the walk
\(W^{(2t-1)}\) followed by the reverse of \(W^{(2t)}\) goes from a vertex of
\(C^{(2t-1)}\) to a vertex of \(C^{(2t)}\). It avoids \(E(C^{(2t-1)})\) and
\(E(C^{(2t)})\): each of the two walks avoids its own cycle by the definition
of a package, and the other cycle because the two arrivals come from different
children. By Lemma~\ref{lem:walk} we obtain a barbell with ends
\(C^{(2t-1)},C^{(2t)}\), whose edges lie in the union of the edge sets of the
two arrivals, and we add it to \(\mathcal B_x\). If \(r\) is odd, we export the
last arrival, \(\pi_x:=(C^{(r)},W^{(r)})\), which is a package at \(v=\hat x\).
Nothing is added to \(\mathcal L_x\).

\paragraph{Case 2: \(x=c_i\) is a cycle node.}
Write \(C:=C_i\) and \(p:=p_i\). The arrivals are packages
\((C^{(1)},W^{(1)}),\dots,(C^{(r)},W^{(r)})\) at distinct vertices
\(a_1,\dots,a_r\) of \(C\), all different from \(p\). None of their edges is an
edge of \(C\), since they lie in the resources of the children, which do not
contain \(c_i\).
\begin{itemize}
\item \(r=0\). Export \(\pi_x:=(C,p)\), where \(p\) denotes the walk of length
  \(0\) at \(p\). This is a package at \(\hat x=p\) with edge set
  \(E(C)\subseteq R_x\).
\item \(r=1\). The walk \(W^{(1)}\) goes from a vertex of \(C^{(1)}\) to
  \(a_1\in V(C)\), and it avoids \(E(C^{(1)})\) and \(E(C)\). By
  Lemma~\ref{lem:walk} we obtain a barbell with ends \(C^{(1)}\) and \(C\),
  whose edges lie in \(E(C^{(1)},W^{(1)})\cup E(C)\). We add it to
  \(\mathcal B_x\) and export nothing.
\item \(r\ge2\). Then \(C\) contains the \(r+1\ge3\) distinct vertices
  \(p,a_1,\dots,a_r\), so its length is \(L\ge3\), and we can write its
  vertices as \(v_0=p,v_1,\dots,v_{L-1}\) in cyclic order. Relabel the arrivals
  so that \(a_j=v_{i_j}\) with \(0<i_1<\cdots<i_r<L\), and let
  \(A_1,\dots,A_r\) be the arcs of Lemma~\ref{lem:arcs}. For
  \(t=1,\dots,\lfloor r/2\rfloor\), the walk
  \[
    W^{(2t-1)}\;A_{2t-1}\;\overline{W^{(2t)}},
  \]
  where the bar denotes reversal, goes from a vertex of \(C^{(2t-1)}\) to a
  vertex of \(C^{(2t)}\). Its edges in \(A_{2t-1}\) are edges of \(C\), which
  is edge-disjoint from \(C^{(2t-1)}\) and \(C^{(2t)}\), and its other edges
  avoid these two cycles as in Case~1. By Lemma~\ref{lem:walk} we obtain a
  barbell with ends \(C^{(2t-1)},C^{(2t)}\), and we add it to
  \(\mathcal B_x\). If \(r\) is odd, we export
  \(\pi_x:=(C^{(r)},W^{(r)}A_r)\), which is a package at \(v_0=p=\hat x\),
  because \(A_r\) ends at \(v_0\) and is edge-disjoint from \(C^{(r)}\).
  Finally we add \(C\) to \(\mathcal L_x\).
\end{itemize}

\paragraph{Checking the properties.}
\emph{\ref{P:partition}.} The cycles \(C_i\) with \(c_i\in T_x\) are the cycles
of the subtrees of the children, together with \(C\) in Case~2. Each cycle of
a child's subtree that was exported by that child is now an end of exactly one
new barbell, or is exported by \(x\). The other cycles of the children keep
their status. In Case~2 the cycle \(C\) is exported (\(r=0\)), an end of one
new barbell (\(r=1\)), or lost (\(r\ge2\)).

\emph{\ref{P:disjoint}.} Every new barbell and the new exported package use
edges of distinct arrivals, together with, in Case~2, edges of \(C\). The
arrivals are pairwise edge-disjoint and edge-disjoint from the old barbells.
In Case~2 the edges of \(C\) belong to \(R_x\) but to no \(R_{y_j}\), so they
are not used by the arrivals or the old barbells. Within Case~2, they are used
only once: by the single barbell when \(r=1\), as the exported cycle when
\(r=0\), and, when \(r\ge2\), through the arcs \(A_1,A_3,\dots\) and \(A_r\),
which are pairwise edge-disjoint by Lemma~\ref{lem:arcs}. All these edges lie
in \(R_x\).

\emph{\ref{P:count}.} A cycle is added to \(\mathcal L_x\) only in Case~2 with
\(r\ge2\), and then at least \(\lfloor r/2\rfloor\ge1\) barbells are added to
\(\mathcal B_x\) at the same time.

This completes the induction.

\paragraph{Conclusion.}
Apply the claim at the root \(\rho\). Since \(T\) spans \(A\), every cycle node
lies in \(T_\rho\). By~\ref{P:partition}, each of the \(k\) cycles is an end of
exactly one barbell of \(\mathcal B_\rho\), or lost, or exported by the root.
Every barbell has exactly two ends, so with \(b:=|\mathcal B_\rho|\),
\[
  k\le 2b+|\mathcal L_\rho|+1\le 3b+1
\]
by~\ref{P:count}. Hence \(b\ge(k-1)/3\), and since \(b\) is an integer,
\(b\ge\lceil(k-1)/3\rceil\). By~\ref{P:disjoint} the barbells are
edge-disjoint, and by~\ref{P:partition} their ends are cycles \(C_i\).

If every \(C_i\) has length at most \(2\), then \(|V(C_i)|\le2\), so Case~2 with
\(r\ge2\) never occurs, because it requires three distinct vertices on \(C\).
Then \(\mathcal L_\rho=\emptyset\) and \(k\le2b+1\), that is,
\(b\ge\lfloor k/2\rfloor\). \qed

\section{The bound is exact}

\begin{theorem}\label{thm:tight}
For every \(m\ge1\) there is a connected simple graph with \(k=3m\)
edge-disjoint triangles that contains at most \(m=\lceil(k-1)/3\rceil\)
edge-disjoint barbells.
\end{theorem}

\begin{proof}
Take a hub vertex \(h\) and \(m\) gadgets. Gadget \(g\) consists of a triangle
\(Z_g=p_ga_gb_g\), a triangle \(X_g\) that meets \(Z_g\) only in the vertex
\(a_g\), a triangle \(Y_g\) that meets \(Z_g\) only in \(b_g\), and the edge
\(hp_g\). All other vertices are new. The graph \(G\) is connected and simple,
and every edge lies in at most one cycle. Its cycles are exactly the \(3m\)
triangles, and they are edge-disjoint.

Consider any family of edge-disjoint barbells. A triangle is an end of at most
one of them, because two barbells with a common end share edges. We claim that
in each gadget at most two of the three triangles are ends. Suppose that
\(X_g\), \(Y_g\) and \(Z_g\) are all ends. Let \(B\) be the barbell with end
\(X_g\), let \(D\) be its other end, and let \(P\) be its path.

In \(G-E(Z_g)\), the vertices \(a_g\), \(b_g\) and \(p_g\) lie in three
different components, and \(V(X_g)\) lies in the component of \(a_g\). If
\(D\ne Z_g\), then \(D\) is \(Y_g\) or a triangle of another gadget, so
\(V(D)\) lies in the component of \(b_g\) or of \(p_g\). The cycles \(X_g\) and
\(D\) have no common vertex, so \(P\) joins two different components of
\(G-E(Z_g)\) and therefore uses an edge of \(Z_g\). But \(Z_g\) is an end of a
barbell different from \(B\), because \(B\) has only the ends \(X_g\) and
\(D\). This contradicts edge-disjointness, so \(D=Z_g\). By the same argument
with \(Y_g\) in place of \(X_g\), the barbell with end \(Y_g\) also has
\(Z_g\) as its other end. These two barbells are different, since a barbell has
only two ends, and both contain \(E(Z_g)\). This is a contradiction.

Hence at most \(2m\) triangles are ends, and since every barbell has two ends
that are triangles, there are at most \(m\) barbells. Finally,
\(\lceil(3m-1)/3\rceil=m\), so Theorem~\ref{thm:main} is exact for this graph.
\end{proof}

\begin{remark}
In this example \(\alpha=3\), so the bound \(k/(2\alpha)=k/6\) of
Corollary~\ref{cor:alpha} is off by a factor of \(2\), while the bound of
Theorem~\ref{thm:main} is exact. The length bound \(\alpha\) is not needed for
Theorem~\ref{thm:main} at all; it enters only when the result is compared with
\(k/(2\alpha)\).
\end{remark}

\end{document}
